\chapter{Catarrh}

\section*{Definition}
Catarrh is excess mucus and swelling in the nose, sinuses, or throat. It can cause a blocked or runny nose, postnasal drip, throat clearing, cough, or a sensation of mucus in the throat.

\section*{When to See a Doctor}

\begin{itemize}
  \item Symptoms persist beyond the typical cold duration, especially if one-sided or with blood-tinged discharge
\end{itemize}

\section*{How It Develops}
Viral infections, allergies, irritants, or inflammation increase mucus production and swelling in the lining of the nose and upper airway. Postnasal drainage can irritate the throat and trigger cough or throat clearing. Persistent or one-sided symptoms may have a structural or less common cause.

\section*{Causes}
\begin{itemize}
 \item Allergic rhinitis hay fever
 \item Upper respiratory viral infections
 \item Sinusitis acute chronic
 \item Gastroesophageal reflux disease
 \item Environmental irritants pollutants smoke
\end{itemize}

\section*{Related Symptoms}
\begin{itemize}
 \item Sneezing runny nose sore throat cough
 \item Postnasal drip headache facial pressure
 \item Reduced sense of smell sleep disturbance snoring daytime fatigue
 \item Cognitive impairment concentration difficulty
 \item Voice changes hoarseness laryngitis
\end{itemize}

\section*{Medical Evaluation}
\begin{itemize}
 \item History focuses on onset, duration, character of symptoms (productive vs dry cough, nasal discharge color), fever pattern, and exposure history.
 \item Physical examination includes vital signs, lung auscultation, nasal inspection, and oropharyngeal examination.
 \item Tests are not usually needed for uncomplicated catarrh. Targeted testing may be considered when another diagnosis is suspected, and allergy testing or lung-function testing may help in selected chronic or recurrent cases.
\end{itemize}

\section*{Treatment Options}
\begin{itemize}
  \item Saline nasal spray or irrigation, avoidance of irritants, and treatment of allergic rhinitis may relieve catarrh.
  \item Persistent one-sided discharge, blood, facial swelling, high fever, or breathing difficulty requires medical assessment.
  \item Antibiotics, nasal corticosteroids, or other treatment should be used when a clinician identifies the underlying cause.
\end{itemize}

\section*{Self-Care and Home Management}
\begin{itemize}
 \item Maintain adequate hydration with water, warm tea, or broth to thin secretions.
 \item A clean humidifier may soothe irritation; avoid very hot steam because of burn risk.
 \item Rest and avoid exposure to smoke, dust, and other respiratory irritants.
 \item Over-the-counter remedies such as saline nasal spray, cough drops, or honey (for adults) may provide symptomatic relief.
 \item Monitor for warning signs including high fever, difficulty breathing, chest pain, or worsening symptoms.
\end{itemize}

\section*{References}
\begin{thebibliography}{99}
\bibitem{catarrh:ref1} Rosenfeld RM, Piccirillo JF, Chandrasekhar SS, et al. ``Clinical Practice Guideline (Update): Adult Sinusitis.'' Otolaryngol Head Neck Surg. 2015;152(2 Suppl):S1--S39.
\bibitem{catarrh:ref2} Fokkens WJ, Lund VJ, Hopkins C, et al. ``European Position Paper on Rhinosinusitis and Nasal Polyps 2020.'' Rhinology. 2020;58(Suppl S29):1--464.
\bibitem{catarrh:ref3} Rosenfeld RM. ``Acute Sinusitis in Adults.'' N Engl J Med. 2016;375(10):962--970.
\bibitem{catarrh:ref4} Wallace DV, Dykewicz MS, Bernstein DI, et al. ``The Diagnosis and Management of Rhinitis: An Updated Practice Parameter.'' J Allergy Clin Immunol. 2008;122(2 Suppl):S1--S84.
\bibitem{catarrh:ref5} Hamilos DL. ``Chronic Rhinosinusitis: Epidemiology and Medical Management.'' J Allergy Clin Immunol. 2011;128(4):693--707.
\end{thebibliography}
